%% Generated by Sphinx.
\def\sphinxdocclass{report}
\documentclass[letterpaper,10pt,english]{sphinxmanual}
\ifdefined\pdfpxdimen
   \let\sphinxpxdimen\pdfpxdimen\else\newdimen\sphinxpxdimen
\fi \sphinxpxdimen=.75bp\relax
\ifdefined\pdfimageresolution
    \pdfimageresolution= \numexpr \dimexpr1in\relax/\sphinxpxdimen\relax
\fi
%% let collapsible pdf bookmarks panel have high depth per default
\PassOptionsToPackage{bookmarksdepth=5}{hyperref}
%% turn off hyperref patch of \index as sphinx.xdy xindy module takes care of
%% suitable \hyperpage mark-up, working around hyperref-xindy incompatibility
\PassOptionsToPackage{hyperindex=false}{hyperref}
%% memoir class requires extra handling
\makeatletter\@ifclassloaded{memoir}
{\ifdefined\memhyperindexfalse\memhyperindexfalse\fi}{}\makeatother

\PassOptionsToPackage{booktabs}{sphinx}
\PassOptionsToPackage{colorrows}{sphinx}

\PassOptionsToPackage{warn}{textcomp}

\catcode`^^^^00a0\active\protected\def^^^^00a0{\leavevmode\nobreak\ }
\usepackage{cmap}
\usepackage{fontspec}
\defaultfontfeatures[\rmfamily,\sffamily,\ttfamily]{}
\usepackage{amsmath,amssymb,amstext}
\usepackage{polyglossia}
\setmainlanguage{english}



\setmainfont{LiberationSerif}[
  UprightFont    = *,
  ItalicFont     = *-Italic,
  BoldFont       = *-Bold,
  BoldItalicFont = *-BoldItalic
]
\setsansfont{LiberationSans}[
  UprightFont    = *,
  ItalicFont     = *-Italic,
  BoldFont       = *-Bold,
  BoldItalicFont = *-BoldItalic
]
\setmonofont{LiberationMono}[
  UprightFont    = *,
  ItalicFont     = *-Italic,
  BoldFont       = *-Bold,
  BoldItalicFont = *-BoldItalic
]



\usepackage[Bjarne]{fncychap}
\usepackage{sphinx}

\fvset{fontsize=auto}
\usepackage{geometry}


% Include hyperref last.
\usepackage{hyperref}
% Fix anchor placement for figures with captions.
\usepackage{hypcap}% it must be loaded after hyperref.
% Set up styles of URL: it should be placed after hyperref.
\urlstyle{same}

\addto\captionsenglish{\renewcommand{\contentsname}{Contents}}

\usepackage{sphinxmessages}
\setcounter{tocdepth}{1}



\title{Pelagos Scientific data on KMaP}
\date{Mar 27, 2026}
\release{0.1}
\author{Pelagos Agreement}
\newcommand{\sphinxlogo}{\vbox{}}
\renewcommand{\releasename}{Release}
\makeindex
\begin{document}

\pagestyle{empty}
\sphinxmaketitle
\pagestyle{plain}
\sphinxtableofcontents
\pagestyle{normal}
\phantomsection\label{\detokenize{index::doc}}


\sphinxAtStartPar
\sphinxstylestrong{User Documentation for the Pelagos Agreement Group}

\begin{sphinxadmonition}{note}{Note:}
\sphinxAtStartPar
This manual is dedicated to Kmap users that are also members of the \sphinxstylestrong{Pelagos Agreement} group
Some features described here may appear diffferently or not be visible to users outside
this group.
\end{sphinxadmonition}

\sphinxAtStartPar
Following The Permanent Secretariat of the Pelagos Agreementahs signed a Memorandun of Understanding in 2024
to publish the results of scientific calls conmsisting in
maps and geospatial data onto the UNEP\sphinxhyphen{}MAP Knowledge Management Platform (KMaP) managed by INFO/RAC;
The KMaP platform is accessible at: \sphinxurl{https://kmap.info-rac.org}

\sphinxAtStartPar
This manual describe how to access and visualize data, metadata and maps in the KMaP platform

\begin{figure}[htbp]
\centering

\noindent\sphinxincludegraphics{{call1_map}.png}
\end{figure}

\sphinxstepscope


\chapter{Welcome to KMaP}
\label{\detokenize{01_welcome:welcome-to-kmap}}\label{\detokenize{01_welcome:welcome-en}}\label{\detokenize{01_welcome::doc}}

\section{What is KMaP?}
\label{\detokenize{01_welcome:what-is-kmap}}
\sphinxAtStartPar
\sphinxstylestrong{KMaP} (Knowledge Management Platform) is the Spatial Data Infrastructure developed by
\sphinxstylestrong{INFO/RAC}, the Regional Activity Centre for Information and Communication of the
\sphinxhref{https://www.unep.org/unepmap/}{Mediterranean Action Plan} (UNEP/MAP).

\sphinxAtStartPar
KMaP is built on \sphinxhref{https://geonode.org/}{GeoNode}, an open\sphinxhyphen{}source web platform for
managing and sharing geospatial information. It acts as a unified access hub to the
environmental and spatial knowledge heritage of the Mediterranean region.

\sphinxAtStartPar
The platform is structured around three hubs:


\begin{savenotes}\sphinxattablestart
\sphinxthistablewithglobalstyle
\centering
\begin{tabular}[t]{\X{20}{100}\X{15}{100}\X{65}{100}}
\sphinxtoprule
\sphinxstyletheadfamily 
\sphinxAtStartPar
Hub
&\sphinxstyletheadfamily 
\sphinxAtStartPar
Access
&\sphinxstyletheadfamily 
\sphinxAtStartPar
Description
\\
\sphinxmidrule
\sphinxtableatstartofbodyhook
\sphinxAtStartPar
\sphinxstylestrong{Data Hub}
&
\sphinxAtStartPar
\sphinxstyleemphasis{Maps} {\color{red}\bfseries{}|mapsbut|}
&
\sphinxAtStartPar
Geospatial data: datasets, interactive maps, dashboards, and geostories.
\\
\sphinxhline
\sphinxAtStartPar
\sphinxstylestrong{Knowledge Hub}
&
\sphinxAtStartPar
\sphinxstyleemphasis{Library} {\color{red}\bfseries{}|librarybut|}
&
\sphinxAtStartPar
Documentary heritage: reports, publications, and reference documents.
\\
\sphinxhline
\sphinxAtStartPar
\sphinxstylestrong{Exchange Hub}
&
\sphinxAtStartPar
\sphinxstyleemphasis{Network} {\color{red}\bfseries{}|networkbut|}
&
\sphinxAtStartPar
Collaboration and knowledge exchange.
\\
\sphinxbottomrule
\end{tabular}
\sphinxtableafterendhook\par
\sphinxattableend\end{savenotes}

\begin{sphinxadmonition}{note}{Note:}
\sphinxAtStartPar
As a member of the \sphinxstylestrong{Pelagos Agreement} group, you have access to specific resources,
metadata editing rights, and group\sphinxhyphen{}level permissions that are not available to
general public users.
\end{sphinxadmonition}


\section{Resource Types}
\label{\detokenize{01_welcome:resource-types}}
\sphinxAtStartPar
KMaP organises all content into five resource types. Each type serves a distinct purpose.


\subsection{Dataset}
\label{\detokenize{01_welcome:dataset}}
\sphinxAtStartPar
A \sphinxstylestrong{Dataset} is a geographic data layer representing real\sphinxhyphen{}world features or measurements.
It is the foundational building block of the platform — the raw material from which maps,
dashboards, and analyses are constructed.

\sphinxAtStartPar
Datasets can be:
\begin{itemize}
\item {} 
\sphinxAtStartPar
\sphinxstylestrong{Vector} – discrete geographic features represented as points, lines, or polygons
(e.g. coastlines, protected area boundaries, monitoring station locations).

\item {} 
\sphinxAtStartPar
\sphinxstylestrong{Raster} – continuous surfaces represented as a grid of cells
(e.g. sea surface temperature, satellite imagery, depth models).

\end{itemize}

\sphinxAtStartPar
What you can do with a Dataset:
\begin{itemize}
\item {} 
\sphinxAtStartPar
View it interactively on a web map

\item {} 
\sphinxAtStartPar
Inspect its attribute table

\item {} 
\sphinxAtStartPar
Download it in standard formats (Shapefile, GeoTIFF, GeoJSON, KML, CSV, …)

\item {} 
\sphinxAtStartPar
Read and edit its metadata

\item {} 
\sphinxAtStartPar
Add it as a layer to a Map

\end{itemize}

\begin{sphinxadmonition}{tip}{Tip:}
\sphinxAtStartPar
Datasets published for the Pelagos Agreement are tagged with the keywords
\sphinxcode{\sphinxupquote{pelagos sanctuary}} or \sphinxcode{\sphinxupquote{pelagos agreement}}. Use these keywords to filter
the catalogue quickly (see {\hyperref[\detokenize{02_find_and_visualize:find-en}]{\sphinxcrossref{\DUrole{std}{\DUrole{std-ref}{How to Find and Visualise Your Data}}}}}).
\end{sphinxadmonition}


\subsection{Document}
\label{\detokenize{01_welcome:document}}
\sphinxAtStartPar
A \sphinxstylestrong{Document} is any uploaded file that complements spatial data: a PDF report,
a field photograph, a spreadsheet, a presentation, or any other reference material.

\sphinxAtStartPar
Documents are fully integrated resources in KMaP — they have their own metadata page,
can be searched and filtered, and can be linked to Datasets or Maps.

\sphinxAtStartPar
Supported file types include: \sphinxcode{\sphinxupquote{.pdf}}, \sphinxcode{\sphinxupquote{.docx}}, \sphinxcode{\sphinxupquote{.xlsx}}, \sphinxcode{\sphinxupquote{.pptx}}, \sphinxcode{\sphinxupquote{.jpg}},
\sphinxcode{\sphinxupquote{.png}}, \sphinxcode{\sphinxupquote{.tif}}, \sphinxcode{\sphinxupquote{.mp4}}, \sphinxcode{\sphinxupquote{.zip}}, and many others.


\subsection{Map}
\label{\detokenize{01_welcome:map}}
\sphinxAtStartPar
A \sphinxstylestrong{Map} is an interactive, multi\sphinxhyphen{}layer web map created by composing one or more
Datasets. Maps are built and viewed using the integrated \sphinxstylestrong{MapStore} map viewer.

\sphinxAtStartPar
Each Map:
\begin{itemize}
\item {} 
\sphinxAtStartPar
Combines multiple Dataset layers with individual styles and visibility settings

\item {} 
\sphinxAtStartPar
Includes a basemap (e.g. OpenStreetMap or satellite imagery)

\item {} 
\sphinxAtStartPar
Supports pan, zoom, feature query, measurement, and print tools

\item {} 
\sphinxAtStartPar
Can be shared via a permanent link or embedded in external websites

\end{itemize}


\subsection{GeoStory}
\label{\detokenize{01_welcome:geostory}}
\sphinxAtStartPar
A \sphinxstylestrong{GeoStory} is a scrollable, multimedia narrative that weaves together text,
interactive maps, images, and video — similar to a story map.

\sphinxAtStartPar
GeoStories are composed of sections:


\begin{savenotes}\sphinxattablestart
\sphinxthistablewithglobalstyle
\centering
\begin{tabular}[t]{\X{25}{100}\X{75}{100}}
\sphinxtoprule
\sphinxstyletheadfamily 
\sphinxAtStartPar
Section type
&\sphinxstyletheadfamily 
\sphinxAtStartPar
Description
\\
\sphinxmidrule
\sphinxtableatstartofbodyhook
\sphinxAtStartPar
\sphinxstylestrong{Title}
&
\sphinxAtStartPar
Full\sphinxhyphen{}screen opening panel with background image or video.
\\
\sphinxhline
\sphinxAtStartPar
\sphinxstylestrong{Banner}
&
\sphinxAtStartPar
Horizontal divider or chapter header with text and/or image.
\\
\sphinxhline
\sphinxAtStartPar
\sphinxstylestrong{Paragraph}
&
\sphinxAtStartPar
Free\sphinxhyphen{}form text and media (image, video, map).
\\
\sphinxhline
\sphinxAtStartPar
\sphinxstylestrong{Immersive}
&
\sphinxAtStartPar
Full\sphinxhyphen{}screen interactive map that the text scrolls over.
\\
\sphinxhline
\sphinxAtStartPar
\sphinxstylestrong{GeoCarousel}
&
\sphinxAtStartPar
Slideshow linked to map locations.
\\
\sphinxhline
\sphinxAtStartPar
\sphinxstylestrong{Web page}
&
\sphinxAtStartPar
Embedded live webpage displayed inline.
\\
\sphinxbottomrule
\end{tabular}
\sphinxtableafterendhook\par
\sphinxattableend\end{savenotes}

\sphinxAtStartPar
GeoStories are ideal for communicating environmental topics to non\sphinxhyphen{}specialist audiences.


\subsection{Dashboard}
\label{\detokenize{01_welcome:dashboard}}
\sphinxAtStartPar
A \sphinxstylestrong{Dashboard} is a configurable panel of interactive widgets — charts, tables, maps,
counters, and text blocks — linked to one or more Datasets.

\sphinxAtStartPar
Available widget types:


\begin{savenotes}\sphinxattablestart
\sphinxthistablewithglobalstyle
\centering
\begin{tabular}[t]{\X{25}{100}\X{75}{100}}
\sphinxtoprule
\sphinxstyletheadfamily 
\sphinxAtStartPar
Widget
&\sphinxstyletheadfamily 
\sphinxAtStartPar
Description
\\
\sphinxmidrule
\sphinxtableatstartofbodyhook
\sphinxAtStartPar
\sphinxstylestrong{Map}
&
\sphinxAtStartPar
Interactive mini\sphinxhyphen{}map displaying dataset layers.
\\
\sphinxhline
\sphinxAtStartPar
\sphinxstylestrong{Chart}
&
\sphinxAtStartPar
Bar, line, pie, or scatter plots derived from dataset attributes.
\\
\sphinxhline
\sphinxAtStartPar
\sphinxstylestrong{Table}
&
\sphinxAtStartPar
Sortable and filterable tabular view of dataset attributes.
\\
\sphinxhline
\sphinxAtStartPar
\sphinxstylestrong{Counter}
&
\sphinxAtStartPar
Large\sphinxhyphen{}format summary statistic (count, sum, average, …).
\\
\sphinxhline
\sphinxAtStartPar
\sphinxstylestrong{Text}
&
\sphinxAtStartPar
Free\sphinxhyphen{}text block for titles or commentary.
\\
\sphinxbottomrule
\end{tabular}
\sphinxtableafterendhook\par
\sphinxattableend\end{savenotes}

\sphinxAtStartPar
Widgets in a Dashboard are interconnected: selecting a region on the map widget
automatically filters all other widgets to display data for that region only.


\section{Interoperable Services}
\label{\detokenize{01_welcome:interoperable-services}}
\sphinxAtStartPar
KMaP exposes its data through standard \sphinxstylestrong{OGC web services}, allowing any compatible
GIS software or web application to access the data directly without downloading files.


\begin{savenotes}\sphinxattablestart
\sphinxthistablewithglobalstyle
\centering
\begin{tabular}[t]{\X{15}{100}\X{15}{100}\X{70}{100}}
\sphinxtoprule
\sphinxstyletheadfamily 
\sphinxAtStartPar
Service
&\sphinxstyletheadfamily 
\sphinxAtStartPar
Standard
&\sphinxstyletheadfamily 
\sphinxAtStartPar
Description
\\
\sphinxmidrule
\sphinxtableatstartofbodyhook
\sphinxAtStartPar
\sphinxstylestrong{WMS}
&
\sphinxAtStartPar
OGC Web Map Service
&
\sphinxAtStartPar
Returns dataset layers as rendered map images. Ideal for visualisation in
desktop GIS (QGIS, ArcGIS) or web applications.
\\
\sphinxhline
\sphinxAtStartPar
\sphinxstylestrong{WFS}
&
\sphinxAtStartPar
OGC Web Feature Service
&
\sphinxAtStartPar
Returns raw vector feature data (geometries + attributes) for download or
analysis. Supports filtering by attribute or bounding box.
\\
\sphinxhline
\sphinxAtStartPar
\sphinxstylestrong{WCS}
&
\sphinxAtStartPar
OGC Web Coverage Service
&
\sphinxAtStartPar
Returns raw raster data (grids) for analysis. Available for raster datasets.
\\
\sphinxhline
\sphinxAtStartPar
\sphinxstylestrong{WMTS}
&
\sphinxAtStartPar
OGC Web Map Tile Service
&
\sphinxAtStartPar
Returns pre\sphinxhyphen{}rendered map tiles for fast display in web maps.
\\
\sphinxhline
\sphinxAtStartPar
\sphinxstylestrong{CSW}
&
\sphinxAtStartPar
OGC Catalogue Service
&
\sphinxAtStartPar
Exposes KMaP’s metadata catalogue for search and harvest by external portals.
\\
\sphinxbottomrule
\end{tabular}
\sphinxtableafterendhook\par
\sphinxattableend\end{savenotes}
\subsubsection*{How to get service URLs}

\sphinxAtStartPar
On any Dataset detail page, click the \sphinxstylestrong{OGC Services} tab (or look for the
\sphinxstyleemphasis{Download / OGC} section). The WMS, WFS, and WCS endpoint URLs are listed there
and can be copied directly into QGIS or any other OGC\sphinxhyphen{}compatible client.

\begin{figure}[htbp]
\centering
\capstart

\noindent\sphinxincludegraphics{{ogc_services}.png}
\caption{The ogc service links in metadata section}\label{\detokenize{01_welcome:id1}}\end{figure}

\begin{sphinxadmonition}{tip}{Tip:}
\sphinxAtStartPar
In QGIS, go to \sphinxstylestrong{Layer → Add Layer → Add WMS/WMTS Layer}, paste the WMS URL,
and browse the list of available layers to add KMaP data directly to your project
without downloading any files.
Consider also the \sphinxhref{https://geonode.org/QGISGeoNodePlugin/}{QGIS Geonode Plugin}
to connect directly with kmap resources.
\end{sphinxadmonition}


\section{Themes Available on KMaP}
\label{\detokenize{01_welcome:themes-available-on-kmap}}
\sphinxAtStartPar
KMaP organises its datasets into the the following categories according with UNEP\sphinxhyphen{}MAP themes:
\begin{itemize}
\item {} 
\sphinxAtStartPar
\sphinxincludegraphics[height=3.5em]{{01pollution}.png} \sphinxstylestrong{Pollution}

\item {} 
\sphinxAtStartPar
\sphinxincludegraphics[height=3.5em]{{02biodiversity}.png} \sphinxstylestrong{Marine Biodiversity}

\item {} 
\sphinxAtStartPar
\sphinxincludegraphics[height=3.5em]{{03sustainability}.png} \sphinxstylestrong{Sustainability and Blue Economy}

\item {} 
\sphinxAtStartPar
\sphinxincludegraphics[height=3.5em]{{04climate_change}.png} \sphinxstylestrong{Climate Change}

\item {} 
\sphinxAtStartPar
\sphinxincludegraphics[height=3.5em]{{05MSP}.png} \sphinxstylestrong{Marine Spatial Planning}

\item {} 
\sphinxAtStartPar
\sphinxincludegraphics[height=3.5em]{{06fishery}.png} \sphinxstylestrong{Fishery and Aquaculture}

\item {} 
\sphinxAtStartPar
\sphinxincludegraphics[height=3.5em]{{07governance}.png} \sphinxstylestrong{Governance}

\end{itemize}

\sphinxAtStartPar
Resources related to the \sphinxhref{https://pelagos-sanctuary.org/scientific-studies/}{Scientific Studies financied by Pelagos Agreement}
are categorized under dtheir respective themes, mainly Marine Biodiversity  and Pollution.

\sphinxstepscope


\chapter{How to Find and Visualise Your Data}
\label{\detokenize{02_find_and_visualize:how-to-find-and-visualise-your-data}}\label{\detokenize{02_find_and_visualize:find-en}}\label{\detokenize{02_find_and_visualize::doc}}
\begin{sphinxcontents}
\sphinxstylecontentstitle{On this page}
\begin{itemize}
\item {} 
\sphinxAtStartPar
\phantomsection\label{\detokenize{02_find_and_visualize:id1}}{\hyperref[\detokenize{02_find_and_visualize:accessing-the-catalogue}]{\sphinxcrossref{Accessing the Catalogue}}}

\item {} 
\sphinxAtStartPar
\phantomsection\label{\detokenize{02_find_and_visualize:id2}}{\hyperref[\detokenize{02_find_and_visualize:filtering-and-searching-the-catalogue}]{\sphinxcrossref{Filtering and Searching the Catalogue}}}
\begin{itemize}
\item {} 
\sphinxAtStartPar
\phantomsection\label{\detokenize{02_find_and_visualize:id3}}{\hyperref[\detokenize{02_find_and_visualize:using-the-search-bar}]{\sphinxcrossref{Using the Search Bar}}}

\item {} 
\sphinxAtStartPar
\phantomsection\label{\detokenize{02_find_and_visualize:id4}}{\hyperref[\detokenize{02_find_and_visualize:using-filters}]{\sphinxcrossref{Using Filters}}}

\end{itemize}

\item {} 
\sphinxAtStartPar
\phantomsection\label{\detokenize{02_find_and_visualize:id5}}{\hyperref[\detokenize{02_find_and_visualize:step-by-step-example-finding-a-pelagos-dataset}]{\sphinxcrossref{Step\sphinxhyphen{}by\sphinxhyphen{}Step Example: Finding a Pelagos Dataset}}}
\begin{itemize}
\item {} 
\sphinxAtStartPar
\phantomsection\label{\detokenize{02_find_and_visualize:id6}}{\hyperref[\detokenize{02_find_and_visualize:step-1-open-the-catalogue}]{\sphinxcrossref{Step 1 – Open the catalogue}}}

\item {} 
\sphinxAtStartPar
\phantomsection\label{\detokenize{02_find_and_visualize:id7}}{\hyperref[\detokenize{02_find_and_visualize:step-2-search-by-keyword}]{\sphinxcrossref{Step 2 – Search by keyword}}}

\item {} 
\sphinxAtStartPar
\phantomsection\label{\detokenize{02_find_and_visualize:id8}}{\hyperref[\detokenize{02_find_and_visualize:step-3-filter-to-datasets-only}]{\sphinxcrossref{Step 3 – Filter to Datasets only}}}

\item {} 
\sphinxAtStartPar
\phantomsection\label{\detokenize{02_find_and_visualize:id9}}{\hyperref[\detokenize{02_find_and_visualize:step-4-select-a-dataset}]{\sphinxcrossref{Step 4 – Select a Dataset}}}

\end{itemize}

\item {} 
\sphinxAtStartPar
\phantomsection\label{\detokenize{02_find_and_visualize:id10}}{\hyperref[\detokenize{02_find_and_visualize:the-dataset-detail-page}]{\sphinxcrossref{The Dataset Detail Page}}}
\begin{itemize}
\item {} 
\sphinxAtStartPar
\phantomsection\label{\detokenize{02_find_and_visualize:id11}}{\hyperref[\detokenize{02_find_and_visualize:thumbnail-and-actions}]{\sphinxcrossref{Thumbnail and Actions}}}

\item {} 
\sphinxAtStartPar
\phantomsection\label{\detokenize{02_find_and_visualize:id12}}{\hyperref[\detokenize{02_find_and_visualize:interactive-preview}]{\sphinxcrossref{Interactive Preview}}}

\item {} 
\sphinxAtStartPar
\phantomsection\label{\detokenize{02_find_and_visualize:id13}}{\hyperref[\detokenize{02_find_and_visualize:metadata-panel}]{\sphinxcrossref{Metadata Panel}}}

\item {} 
\sphinxAtStartPar
\phantomsection\label{\detokenize{02_find_and_visualize:id14}}{\hyperref[\detokenize{02_find_and_visualize:full-metadata-view}]{\sphinxcrossref{Full Metadata View}}}

\item {} 
\sphinxAtStartPar
\phantomsection\label{\detokenize{02_find_and_visualize:id15}}{\hyperref[\detokenize{02_find_and_visualize:downloading-a-dataset}]{\sphinxcrossref{Downloading a Dataset}}}

\end{itemize}

\end{itemize}
\end{sphinxcontents}

\sphinxAtStartPar
This section explains how to discover resources in KMaP, apply filters to narrow your
search, and explore the detail page and preview of a Dataset.


\section{Accessing the Catalogue}
\label{\detokenize{02_find_and_visualize:accessing-the-catalogue}}
\sphinxAtStartPar
The KMaP catalogue lists all public resources and all resources shared with your group.
There are two main ways to reach it:
\begin{enumerate}
\sphinxsetlistlabels{\arabic}{enumi}{enumii}{}{.}%
\item {} 
\sphinxAtStartPar
\sphinxstylestrong{From the navigation bar} – click \sphinxstylestrong{Maps} (for spatial resources) or \sphinxstylestrong{Library}
(for documents and publications) in the top menu.

\item {} 
\sphinxAtStartPar
\sphinxstylestrong{Directly via URL} – navigate to \sphinxurl{https://kmap.info-rac.org/resources/}

\end{enumerate}

\sphinxAtStartPar
The catalogue page displays resources as cards, each showing a thumbnail, title, resource
type, and summary metadata.

\begin{sphinxadmonition}{tip}{Tip:}
\sphinxAtStartPar
Make sure you are \sphinxstylestrong{logged in} to see resources that are shared exclusively with the
Pelagos Agreement group. Resources restricted to your group are not visible to
unauthenticated users.
\end{sphinxadmonition}


\section{Filtering and Searching the Catalogue}
\label{\detokenize{02_find_and_visualize:filtering-and-searching-the-catalogue}}

\subsection{Using the Search Bar}
\label{\detokenize{02_find_and_visualize:using-the-search-bar}}
\sphinxAtStartPar
The search bar at the top of the catalogue page accepts free\sphinxhyphen{}text queries. It searches
across titles, abstracts, and keyword fields.

\sphinxAtStartPar
To find Pelagos\sphinxhyphen{}related content:
\begin{enumerate}
\sphinxsetlistlabels{\arabic}{enumi}{enumii}{}{.}%
\item {} 
\sphinxAtStartPar
Type \sphinxcode{\sphinxupquote{pelagos sanctuary}} in the search bar and press \sphinxstylestrong{Enter}.

\item {} 
\sphinxAtStartPar
Alternatively, search for \sphinxcode{\sphinxupquote{pelagos agreement}} to retrieve resources tagged
with that keyword.

\end{enumerate}

\begin{sphinxadmonition}{note}{Note:}
\sphinxAtStartPar
Both \sphinxcode{\sphinxupquote{pelagos sanctuary}} and \sphinxcode{\sphinxupquote{pelagos agreement}} are controlled keywords used by
the Pelagos Agreement group when publishing resources on KMaP. Using these exact terms
will return the most relevant results.
\end{sphinxadmonition}


\subsection{Using Filters}
\label{\detokenize{02_find_and_visualize:using-filters}}
\sphinxAtStartPar
On the left side of the catalogue page you will find a filter panel. Filters can be
combined freely and are applied immediately.


\begin{savenotes}\sphinxattablestart
\sphinxthistablewithglobalstyle
\centering
\begin{tabular}[t]{\X{25}{100}\X{75}{100}}
\sphinxtoprule
\sphinxstyletheadfamily 
\sphinxAtStartPar
Filter
&\sphinxstyletheadfamily 
\sphinxAtStartPar
How to use it
\\
\sphinxmidrule
\sphinxtableatstartofbodyhook
\sphinxAtStartPar
\sphinxstylestrong{Resource type}
&
\sphinxAtStartPar
Select \sphinxstyleemphasis{Dataset}, \sphinxstyleemphasis{Document}, \sphinxstyleemphasis{Map}, \sphinxstyleemphasis{GeoStory}, or \sphinxstyleemphasis{Dashboard} to restrict
results to one type.
\\
\sphinxhline
\sphinxAtStartPar
\sphinxstylestrong{Keywords}
&
\sphinxAtStartPar
Type or select a keyword (e.g. \sphinxcode{\sphinxupquote{pelagos sanctuary}}) to filter by thematic tag.
\\
\sphinxhline
\sphinxAtStartPar
\sphinxstylestrong{Category / Theme}
&
\sphinxAtStartPar
Select a thematic category such as \sphinxstyleemphasis{Marine Biodiversity} or \sphinxstyleemphasis{Governance}.
\\
\sphinxhline
\sphinxAtStartPar
\sphinxstylestrong{Date}
&
\sphinxAtStartPar
Filter by publication or modification date range.
\\
\sphinxhline
\sphinxAtStartPar
\sphinxstylestrong{Owner}
&
\sphinxAtStartPar
Filter by the user or organisation that published the resource.
\\
\sphinxhline
\sphinxAtStartPar
\sphinxstylestrong{Extent}
&
\sphinxAtStartPar
Draw a bounding box on the mini\sphinxhyphen{}map to filter by geographic area.
\\
\sphinxbottomrule
\end{tabular}
\sphinxtableafterendhook\par
\sphinxattableend\end{savenotes}

\begin{sphinxadmonition}{tip}{Tip:}
\sphinxAtStartPar
To find all Pelagos\sphinxhyphen{}related datasets quickly:
\begin{enumerate}
\sphinxsetlistlabels{\arabic}{enumi}{enumii}{}{.}%
\item {} 
\sphinxAtStartPar
Set \sphinxstylestrong{Resource type} → \sphinxstyleemphasis{Dataset}

\item {} 
\sphinxAtStartPar
In the \sphinxstylestrong{Keywords} filter, type \sphinxcode{\sphinxupquote{pelagos sanctuary}}

\item {} 
\sphinxAtStartPar
The catalogue will update to show only datasets tagged with that keyword.

\end{enumerate}
\end{sphinxadmonition}


\section{Step\sphinxhyphen{}by\sphinxhyphen{}Step Example: Finding a Pelagos Dataset}
\label{\detokenize{02_find_and_visualize:step-by-step-example-finding-a-pelagos-dataset}}
\sphinxAtStartPar
The following example walks through finding, previewing, and reading the metadata of a
Dataset relevant to the Pelagos Agreement.


\subsection{Step 1 – Open the catalogue}
\label{\detokenize{02_find_and_visualize:step-1-open-the-catalogue}}
\sphinxAtStartPar
Go to \sphinxurl{https://kmap.info-rac.org} and click \sphinxstylestrong{Maps} in the top navigation bar.
This opens the Data Hub catalogue, filtered to spatial resource types.


\subsection{Step 2 – Search by keyword}
\label{\detokenize{02_find_and_visualize:step-2-search-by-keyword}}
\sphinxAtStartPar
In the search bar, type:

\begin{sphinxVerbatim}[commandchars=\\\{\}]
\PYG{n}{pelagos} \PYG{n}{sanctuary}
\end{sphinxVerbatim}

\sphinxAtStartPar
Press \sphinxstylestrong{Enter}. The results panel updates to show resources matching that keyword.


\subsection{Step 3 – Filter to Datasets only}
\label{\detokenize{02_find_and_visualize:step-3-filter-to-datasets-only}}
\sphinxAtStartPar
In the left filter panel, under \sphinxstylestrong{Resource type}, tick \sphinxstylestrong{Dataset}.
The list now shows only datasets tagged with \sphinxstyleemphasis{pelagos sanctuary}.


\subsection{Step 4 – Select a Dataset}
\label{\detokenize{02_find_and_visualize:step-4-select-a-dataset}}
\sphinxAtStartPar
Click on a dataset card to open its \sphinxstylestrong{detail page}.


\section{The Dataset Detail Page}
\label{\detokenize{02_find_and_visualize:the-dataset-detail-page}}
\sphinxAtStartPar
The detail page is the central information page for any resource. It is divided into
several sections.


\subsection{Thumbnail and Actions}
\label{\detokenize{02_find_and_visualize:thumbnail-and-actions}}
\sphinxAtStartPar
At the top of the page you will see:
\begin{itemize}
\item {} 
\sphinxAtStartPar
A \sphinxstylestrong{map thumbnail} or preview image of the dataset.

\item {} 
\sphinxAtStartPar
Action buttons:
\begin{itemize}
\item {} 
\sphinxAtStartPar
\sphinxstylestrong{View map} – opens the dataset in the interactive map viewer.

\item {} 
\sphinxAtStartPar
\sphinxstylestrong{Download} – downloads the dataset in your chosen format.

\item {} 
\sphinxAtStartPar
\sphinxstylestrong{Edit} (if you have permission) – opens the metadata editor.

\item {} 
\sphinxAtStartPar
\sphinxstylestrong{Share} – copies the permalink to the resource.

\end{itemize}

\end{itemize}


\subsection{Interactive Preview}
\label{\detokenize{02_find_and_visualize:interactive-preview}}
\sphinxAtStartPar
Below the thumbnail, the \sphinxstylestrong{map preview panel} displays the dataset rendered on a
basemap. You can:
\begin{itemize}
\item {} 
\sphinxAtStartPar
Pan and zoom using mouse or touch gestures.

\item {} 
\sphinxAtStartPar
Click on a feature to open a \sphinxstylestrong{pop\sphinxhyphen{}up} showing its attribute values.

\item {} 
\sphinxAtStartPar
Use the layer panel (☰ icon) to toggle visibility or inspect the legend.

\end{itemize}

\begin{sphinxadmonition}{note}{Note:}
\sphinxAtStartPar
The preview is a live, interactive map — not a static image. You can explore the
data spatially before deciding whether to download or use it.
\end{sphinxadmonition}


\subsection{Metadata Panel}
\label{\detokenize{02_find_and_visualize:metadata-panel}}
\sphinxAtStartPar
Below the preview, the \sphinxstylestrong{metadata panel} contains detailed information about the
dataset. The key fields are:


\begin{savenotes}\sphinxattablestart
\sphinxthistablewithglobalstyle
\centering
\begin{tabular}[t]{\X{30}{100}\X{70}{100}}
\sphinxtoprule
\sphinxstyletheadfamily 
\sphinxAtStartPar
Field
&\sphinxstyletheadfamily 
\sphinxAtStartPar
Description
\\
\sphinxmidrule
\sphinxtableatstartofbodyhook
\sphinxAtStartPar
\sphinxstylestrong{Title}
&
\sphinxAtStartPar
The full name of the dataset.
\\
\sphinxhline
\sphinxAtStartPar
\sphinxstylestrong{Abstract}
&
\sphinxAtStartPar
A plain\sphinxhyphen{}language description of what the dataset contains and its purpose.
\\
\sphinxhline
\sphinxAtStartPar
\sphinxstylestrong{Keywords}
&
\sphinxAtStartPar
Thematic and geographic tags (look for \sphinxcode{\sphinxupquote{pelagos sanctuary}} or
\sphinxcode{\sphinxupquote{pelagos agreement}} here).
\\
\sphinxhline
\sphinxAtStartPar
\sphinxstylestrong{Category}
&
\sphinxAtStartPar
The INSPIRE or thematic category (e.g. \sphinxstyleemphasis{Marine Biodiversity}).
\\
\sphinxhline
\sphinxAtStartPar
\sphinxstylestrong{Date}
&
\sphinxAtStartPar
Publication date and last modification date.
\\
\sphinxhline
\sphinxAtStartPar
\sphinxstylestrong{Spatial extent}
&
\sphinxAtStartPar
Bounding box or geographic region covered by the data.
\\
\sphinxhline
\sphinxAtStartPar
\sphinxstylestrong{Coordinate Reference System (CRS)}
&
\sphinxAtStartPar
The projection used (e.g. WGS 84 / EPSG:4326).
\\
\sphinxhline
\sphinxAtStartPar
\sphinxstylestrong{Owner}
&
\sphinxAtStartPar
The user or organisation responsible for the dataset.
\\
\sphinxhline
\sphinxAtStartPar
\sphinxstylestrong{Licence}
&
\sphinxAtStartPar
The usage licence (e.g. Creative Commons, ODbL, …).
\\
\sphinxhline
\sphinxAtStartPar
\sphinxstylestrong{Point of contact}
&
\sphinxAtStartPar
Who to contact for questions about the data.
\\
\sphinxhline
\sphinxAtStartPar
\sphinxstylestrong{OGC Services}
&
\sphinxAtStartPar
WMS, WFS, and WCS endpoint URLs for direct integration in GIS software.
\\
\sphinxhline
\sphinxAtStartPar
\sphinxstylestrong{Related resources}
&
\sphinxAtStartPar
Links to Maps, Documents, or other Datasets connected to this resource.
\\
\sphinxbottomrule
\end{tabular}
\sphinxtableafterendhook\par
\sphinxattableend\end{savenotes}


\subsection{Full Metadata View}
\label{\detokenize{02_find_and_visualize:full-metadata-view}}
\sphinxAtStartPar
To see the complete metadata record, click the \sphinxstylestrong{Metadata} tab (or the \sphinxstyleemphasis{View full
metadata} link near the bottom of the detail page). This view displays all metadata
fields in their full form, including ISO 19115 / INSPIRE\sphinxhyphen{}compliant fields.


\subsection{Downloading a Dataset}
\label{\detokenize{02_find_and_visualize:downloading-a-dataset}}
\sphinxAtStartPar
Click the \sphinxstylestrong{Download} button on the detail page. A format selector appears. Common
options include:
\begin{itemize}
\item {} 
\sphinxAtStartPar
\sphinxstylestrong{ESRI Shapefile} – for use in QGIS, ArcGIS, and most desktop GIS.

\item {} 
\sphinxAtStartPar
\sphinxstylestrong{GeoJSON} – for web mapping and lightweight GIS use.

\item {} 
\sphinxAtStartPar
\sphinxstylestrong{KML/KMZ} – for Google Earth.

\item {} 
\sphinxAtStartPar
\sphinxstylestrong{GeoPackage} – a single\sphinxhyphen{}file format supported by QGIS.

\item {} 
\sphinxAtStartPar
\sphinxstylestrong{GeoTIFF} – for raster datasets.

\item {} 
\sphinxAtStartPar
\sphinxstylestrong{CSV} – attribute table only (no geometry).

\end{itemize}

\sphinxAtStartPar
Select your preferred format and click \sphinxstylestrong{Download} to save the file.

\sphinxstepscope


\chapter{Managing Metadata}
\label{\detokenize{03_managing_metadata:managing-metadata}}\label{\detokenize{03_managing_metadata:metadata-en}}\label{\detokenize{03_managing_metadata::doc}}
\begin{sphinxcontents}
\sphinxstylecontentstitle{On this page}
\begin{itemize}
\item {} 
\sphinxAtStartPar
\phantomsection\label{\detokenize{03_managing_metadata:id1}}{\hyperref[\detokenize{03_managing_metadata:who-can-edit-metadata}]{\sphinxcrossref{Who Can Edit Metadata?}}}

\item {} 
\sphinxAtStartPar
\phantomsection\label{\detokenize{03_managing_metadata:id2}}{\hyperref[\detokenize{03_managing_metadata:opening-the-metadata-editor}]{\sphinxcrossref{Opening the Metadata Editor}}}

\item {} 
\sphinxAtStartPar
\phantomsection\label{\detokenize{03_managing_metadata:id3}}{\hyperref[\detokenize{03_managing_metadata:overview-of-the-metadata-editor}]{\sphinxcrossref{Overview of the Metadata Editor}}}

\item {} 
\sphinxAtStartPar
\phantomsection\label{\detokenize{03_managing_metadata:id4}}{\hyperref[\detokenize{03_managing_metadata:editing-key-metadata-fields}]{\sphinxcrossref{Editing Key Metadata Fields}}}
\begin{itemize}
\item {} 
\sphinxAtStartPar
\phantomsection\label{\detokenize{03_managing_metadata:id5}}{\hyperref[\detokenize{03_managing_metadata:title}]{\sphinxcrossref{Title}}}

\item {} 
\sphinxAtStartPar
\phantomsection\label{\detokenize{03_managing_metadata:id6}}{\hyperref[\detokenize{03_managing_metadata:abstract}]{\sphinxcrossref{Abstract}}}

\item {} 
\sphinxAtStartPar
\phantomsection\label{\detokenize{03_managing_metadata:id7}}{\hyperref[\detokenize{03_managing_metadata:keywords}]{\sphinxcrossref{Keywords}}}

\item {} 
\sphinxAtStartPar
\phantomsection\label{\detokenize{03_managing_metadata:id8}}{\hyperref[\detokenize{03_managing_metadata:category}]{\sphinxcrossref{Category}}}

\item {} 
\sphinxAtStartPar
\phantomsection\label{\detokenize{03_managing_metadata:id9}}{\hyperref[\detokenize{03_managing_metadata:spatial-extent}]{\sphinxcrossref{Spatial Extent}}}

\item {} 
\sphinxAtStartPar
\phantomsection\label{\detokenize{03_managing_metadata:id10}}{\hyperref[\detokenize{03_managing_metadata:date-fields}]{\sphinxcrossref{Date Fields}}}

\item {} 
\sphinxAtStartPar
\phantomsection\label{\detokenize{03_managing_metadata:id11}}{\hyperref[\detokenize{03_managing_metadata:licence}]{\sphinxcrossref{Licence}}}

\item {} 
\sphinxAtStartPar
\phantomsection\label{\detokenize{03_managing_metadata:id12}}{\hyperref[\detokenize{03_managing_metadata:point-of-contact}]{\sphinxcrossref{Point of Contact}}}

\item {} 
\sphinxAtStartPar
\phantomsection\label{\detokenize{03_managing_metadata:id13}}{\hyperref[\detokenize{03_managing_metadata:thumbnail}]{\sphinxcrossref{Thumbnail}}}

\end{itemize}

\item {} 
\sphinxAtStartPar
\phantomsection\label{\detokenize{03_managing_metadata:id14}}{\hyperref[\detokenize{03_managing_metadata:saving-your-changes}]{\sphinxcrossref{Saving Your Changes}}}

\item {} 
\sphinxAtStartPar
\phantomsection\label{\detokenize{03_managing_metadata:id15}}{\hyperref[\detokenize{03_managing_metadata:metadata-completeness-checklist}]{\sphinxcrossref{Metadata Completeness Checklist}}}

\end{itemize}
\end{sphinxcontents}

\sphinxAtStartPar
Good metadata is what makes a resource discoverable and trustworthy. This section
explains how to view and edit the metadata of a Dataset or Document in KMaP.

\begin{sphinxadmonition}{note}{Note:}
\sphinxAtStartPar
You can only edit the metadata of resources that you \sphinxstylestrong{own} or for which you have
been granted \sphinxstylestrong{edit permissions}. Members of the Pelagos Agreement group may have
edit rights on resources shared within the group — check with your group administrator
if you are unsure.
\end{sphinxadmonition}


\section{Who Can Edit Metadata?}
\label{\detokenize{03_managing_metadata:who-can-edit-metadata}}
\sphinxAtStartPar
KMaP has a permission model with several levels:


\begin{savenotes}\sphinxattablestart
\sphinxthistablewithglobalstyle
\centering
\begin{tabular}[t]{\X{25}{100}\X{75}{100}}
\sphinxtoprule
\sphinxstyletheadfamily 
\sphinxAtStartPar
Role
&\sphinxstyletheadfamily 
\sphinxAtStartPar
Metadata editing rights
\\
\sphinxmidrule
\sphinxtableatstartofbodyhook
\sphinxAtStartPar
\sphinxstylestrong{Resource owner}
&
\sphinxAtStartPar
Full edit rights on all metadata fields.
\\
\sphinxhline
\sphinxAtStartPar
\sphinxstylestrong{Group manager}
&
\sphinxAtStartPar
Can edit resources owned by any member of their group.
\\
\sphinxhline
\sphinxAtStartPar
\sphinxstylestrong{Group member (editor)}
&
\sphinxAtStartPar
Can edit resources explicitly shared with the group with \sphinxstyleemphasis{edit} permission.
\\
\sphinxhline
\sphinxAtStartPar
\sphinxstylestrong{Viewer / public}
&
\sphinxAtStartPar
Read\sphinxhyphen{}only access. Cannot edit metadata.
\\
\sphinxbottomrule
\end{tabular}
\sphinxtableafterendhook\par
\sphinxattableend\end{savenotes}

\sphinxAtStartPar
If you see an \sphinxstylestrong{Edit} button on the resource detail page, you have the right to edit
that resource’s metadata.


\section{Opening the Metadata Editor}
\label{\detokenize{03_managing_metadata:opening-the-metadata-editor}}
\sphinxAtStartPar
There are two ways to open the metadata editor:

\sphinxAtStartPar
\sphinxstylestrong{From the detail page:}
\begin{enumerate}
\sphinxsetlistlabels{\arabic}{enumi}{enumii}{}{.}%
\item {} 
\sphinxAtStartPar
Navigate to the resource detail page (see {\hyperref[\detokenize{02_find_and_visualize:find-en}]{\sphinxcrossref{\DUrole{std}{\DUrole{std-ref}{How to Find and Visualise Your Data}}}}} for how to find a resource).

\item {} 
\sphinxAtStartPar
Click the \sphinxstylestrong{Edit} button (pencil icon) in the top\sphinxhyphen{}right action bar.

\item {} 
\sphinxAtStartPar
A dropdown menu appears. Select \sphinxstylestrong{Edit metadata}.

\end{enumerate}

\sphinxAtStartPar
\sphinxstylestrong{From the catalogue:}
\begin{enumerate}
\sphinxsetlistlabels{\arabic}{enumi}{enumii}{}{.}%
\item {} 
\sphinxAtStartPar
In the catalogue, hover over a resource card.

\item {} 
\sphinxAtStartPar
Click the \sphinxstylestrong{⋮} (more options) menu that appears.

\item {} 
\sphinxAtStartPar
Select \sphinxstylestrong{Edit metadata}.

\end{enumerate}

\sphinxAtStartPar
The metadata editor opens as a multi\sphinxhyphen{}tab form.


\section{Overview of the Metadata Editor}
\label{\detokenize{03_managing_metadata:overview-of-the-metadata-editor}}
\sphinxAtStartPar
The metadata editor is organised into tabs. The most important tabs are:


\begin{savenotes}\sphinxattablestart
\sphinxthistablewithglobalstyle
\centering
\begin{tabular}[t]{\X{25}{100}\X{75}{100}}
\sphinxtoprule
\sphinxstyletheadfamily 
\sphinxAtStartPar
Tab
&\sphinxstyletheadfamily 
\sphinxAtStartPar
Content
\\
\sphinxmidrule
\sphinxtableatstartofbodyhook
\sphinxAtStartPar
\sphinxstylestrong{Basic info}
&
\sphinxAtStartPar
Title, abstract, language, category, keywords, point of contact.
\\
\sphinxhline
\sphinxAtStartPar
\sphinxstylestrong{Location and licences}
&
\sphinxAtStartPar
Spatial extent, CRS, date fields, licence, attribution.
\\
\sphinxhline
\sphinxAtStartPar
\sphinxstylestrong{Optional metadata}
&
\sphinxAtStartPar
Additional ISO 19115 fields (edition, purpose, supplemental information, …).
\\
\sphinxhline
\sphinxAtStartPar
\sphinxstylestrong{Settings}
&
\sphinxAtStartPar
Permissions, featured status, advertised status.
\\
\sphinxbottomrule
\end{tabular}
\sphinxtableafterendhook\par
\sphinxattableend\end{savenotes}


\section{Editing Key Metadata Fields}
\label{\detokenize{03_managing_metadata:editing-key-metadata-fields}}

\subsection{Title}
\label{\detokenize{03_managing_metadata:title}}
\sphinxAtStartPar
The \sphinxstylestrong{title} is the primary way users find your resource. It should be:
\begin{itemize}
\item {} 
\sphinxAtStartPar
Descriptive and specific (avoid generic titles like “Data” or “Map layer”).

\item {} 
\sphinxAtStartPar
Consistent with the naming conventions used by the Pelagos Agreement group.

\end{itemize}

\sphinxAtStartPar
To edit: click on the \sphinxstylestrong{Title} field in the \sphinxstyleemphasis{Basic info} tab and type your changes.


\subsection{Abstract}
\label{\detokenize{03_managing_metadata:abstract}}
\sphinxAtStartPar
The \sphinxstylestrong{abstract} is a plain\sphinxhyphen{}language description of the resource. A good abstract answers:
\begin{itemize}
\item {} 
\sphinxAtStartPar
What does this resource contain?

\item {} 
\sphinxAtStartPar
What geographic area does it cover?

\item {} 
\sphinxAtStartPar
What time period does it represent?

\item {} 
\sphinxAtStartPar
What is it used for?

\item {} 
\sphinxAtStartPar
What are its known limitations?

\end{itemize}

\sphinxAtStartPar
The abstract field supports basic formatting. Aim for at least two or three sentences.


\subsection{Keywords}
\label{\detokenize{03_managing_metadata:keywords}}
\sphinxAtStartPar
Keywords are the primary mechanism for filtering and discovering resources in KMaP.

\begin{sphinxadmonition}{important}{Important:}
\sphinxAtStartPar
All resources related to the Pelagos Agreement \sphinxstylestrong{must include} at least one of the
following keywords:
\begin{itemize}
\item {} 
\sphinxAtStartPar
\sphinxcode{\sphinxupquote{pelagos sanctuary}}

\item {} 
\sphinxAtStartPar
\sphinxcode{\sphinxupquote{pelagos agreement}}

\end{itemize}

\sphinxAtStartPar
Without these keywords, your resource will not appear when other group members
filter the catalogue by these terms.
\end{sphinxadmonition}

\sphinxAtStartPar
To add a keyword:
\begin{enumerate}
\sphinxsetlistlabels{\arabic}{enumi}{enumii}{}{.}%
\item {} 
\sphinxAtStartPar
In the \sphinxstylestrong{Keywords} field, start typing the keyword.

\item {} 
\sphinxAtStartPar
If the keyword already exists in the system, it will appear in a dropdown — select it.

\item {} 
\sphinxAtStartPar
If it is new, type the full keyword and press \sphinxstylestrong{Enter} or click \sphinxstylestrong{Add}.

\item {} 
\sphinxAtStartPar
To remove a keyword, click the \sphinxstylestrong{×} next to it.

\end{enumerate}


\subsection{Category}
\label{\detokenize{03_managing_metadata:category}}
\sphinxAtStartPar
The \sphinxstylestrong{category} assigns the resource to a thematic classification.
Select the most appropriate category from the dropdown:
\begin{itemize}
\item {} 
\sphinxAtStartPar
Fishery and Aquaculture

\item {} 
\sphinxAtStartPar
Marine Biodiversity

\item {} 
\sphinxAtStartPar
Pollution

\item {} 
\sphinxAtStartPar
Climate Change

\item {} 
\sphinxAtStartPar
Marine Spatial Planning

\item {} 
\sphinxAtStartPar
Sustainability and Blue Economy

\item {} 
\sphinxAtStartPar
Governance

\end{itemize}

\sphinxAtStartPar
A resource can belong to only one primary category. Use keywords for additional
thematic tagging.


\subsection{Spatial Extent}
\label{\detokenize{03_managing_metadata:spatial-extent}}
\sphinxAtStartPar
The \sphinxstylestrong{spatial extent} defines the geographic bounding box of the resource. For datasets,
this is usually set automatically when the data is uploaded. You can refine it manually:
\begin{enumerate}
\sphinxsetlistlabels{\arabic}{enumi}{enumii}{}{.}%
\item {} 
\sphinxAtStartPar
Go to the \sphinxstylestrong{Location and licences} tab.

\item {} 
\sphinxAtStartPar
In the \sphinxstyleemphasis{Spatial extent} section, use the map widget to draw or adjust the bounding box,
or enter coordinates (min/max longitude and latitude) manually.

\end{enumerate}

\begin{sphinxadmonition}{tip}{Tip:}
\sphinxAtStartPar
For Pelagos\sphinxhyphen{}related resources, the spatial extent should cover at least the Pelagos
Sanctuary area (approximately between France, Italy, and Monaco in the Ligurian Sea and
adjacent western Mediterranean).
\end{sphinxadmonition}


\subsection{Date Fields}
\label{\detokenize{03_managing_metadata:date-fields}}
\sphinxAtStartPar
Three date fields are available:


\begin{savenotes}\sphinxattablestart
\sphinxthistablewithglobalstyle
\centering
\begin{tabular}[t]{\X{30}{100}\X{70}{100}}
\sphinxtoprule
\sphinxstyletheadfamily 
\sphinxAtStartPar
Field
&\sphinxstyletheadfamily 
\sphinxAtStartPar
Description
\\
\sphinxmidrule
\sphinxtableatstartofbodyhook
\sphinxAtStartPar
\sphinxstylestrong{Date} (creation date)
&
\sphinxAtStartPar
When the resource was originally created or the data was collected.
\\
\sphinxhline
\sphinxAtStartPar
\sphinxstylestrong{Publication date}
&
\sphinxAtStartPar
When the resource was published on KMaP.
\\
\sphinxhline
\sphinxAtStartPar
\sphinxstylestrong{Revision date}
&
\sphinxAtStartPar
When the resource was last updated.
\\
\sphinxbottomrule
\end{tabular}
\sphinxtableafterendhook\par
\sphinxattableend\end{savenotes}

\sphinxAtStartPar
Keep these dates accurate — they are used to filter and sort resources in the catalogue.


\subsection{Licence}
\label{\detokenize{03_managing_metadata:licence}}
\sphinxAtStartPar
The \sphinxstylestrong{licence} defines how others may use the resource.
\begin{enumerate}
\sphinxsetlistlabels{\arabic}{enumi}{enumii}{}{.}%
\item {} 
\sphinxAtStartPar
Go to the \sphinxstylestrong{Location and licences} tab.

\item {} 
\sphinxAtStartPar
Select a licence from the dropdown. Common options include:
\begin{itemize}
\item {} 
\sphinxAtStartPar
\sphinxstyleemphasis{Creative Commons Attribution (CC BY)} – free reuse with attribution.

\item {} 
\sphinxAtStartPar
\sphinxstyleemphasis{Creative Commons Attribution ShareAlike (CC BY\sphinxhyphen{}SA)} – reuse with attribution and
same licence.

\item {} 
\sphinxAtStartPar
\sphinxstyleemphasis{Open Database Licence (ODbL)} – for datasets.

\item {} 
\sphinxAtStartPar
\sphinxstyleemphasis{No restrictions} – public domain.

\end{itemize}

\end{enumerate}

\begin{sphinxadmonition}{note}{Note:}
\sphinxAtStartPar
If you are unsure which licence to apply, consult your group administrator. Applying
the wrong licence may restrict the reuse of your data inappropriately.
\end{sphinxadmonition}


\subsection{Point of Contact}
\label{\detokenize{03_managing_metadata:point-of-contact}}
\sphinxAtStartPar
The \sphinxstylestrong{point of contact} identifies who to contact for questions about the resource.
It can be a person or an organisation.
\begin{enumerate}
\sphinxsetlistlabels{\arabic}{enumi}{enumii}{}{.}%
\item {} 
\sphinxAtStartPar
In the \sphinxstyleemphasis{Basic info} tab, find the \sphinxstylestrong{Point of contact} field.

\item {} 
\sphinxAtStartPar
Start typing a name to search registered KMaP users, or enter an organisation name.

\end{enumerate}


\subsection{Thumbnail}
\label{\detokenize{03_managing_metadata:thumbnail}}
\sphinxAtStartPar
The \sphinxstylestrong{thumbnail} is the preview image displayed in catalogue cards. It is generated
automatically for Datasets and Maps, but can be replaced manually:
\begin{enumerate}
\sphinxsetlistlabels{\arabic}{enumi}{enumii}{}{.}%
\item {} 
\sphinxAtStartPar
Click \sphinxstylestrong{Edit} on the detail page.

\item {} 
\sphinxAtStartPar
Select \sphinxstylestrong{Edit thumbnail}.

\item {} 
\sphinxAtStartPar
Upload a new image (recommended: 600 × 400 px, JPEG or PNG).

\end{enumerate}


\section{Saving Your Changes}
\label{\detokenize{03_managing_metadata:saving-your-changes}}
\sphinxAtStartPar
When you have finished editing, scroll to the bottom of the editor and click
\sphinxstylestrong{Save} (or \sphinxstylestrong{Update} — the label may vary). A confirmation message will appear.

\begin{sphinxadmonition}{warning}{Warning:}
\sphinxAtStartPar
Metadata changes are applied immediately and are visible to all users with access to
the resource. There is no draft or staging mode.
\end{sphinxadmonition}

\sphinxAtStartPar
After saving, you are returned to the resource detail page where you can verify that
all fields have been updated correctly.


\section{Metadata Completeness Checklist}
\label{\detokenize{03_managing_metadata:metadata-completeness-checklist}}
\sphinxAtStartPar
Use the following checklist before publishing or sharing a resource:


\begin{savenotes}\sphinxattablestart
\sphinxthistablewithglobalstyle
\centering
\begin{tabular}[t]{\X{10}{100}\X{90}{100}}
\sphinxtoprule
\sphinxstyletheadfamily 
\sphinxAtStartPar
\(\checkmark\)
&\sphinxstyletheadfamily 
\sphinxAtStartPar
Item
\\
\sphinxmidrule
\sphinxtableatstartofbodyhook
\sphinxAtStartPar
☐
&
\sphinxAtStartPar
Title is descriptive and specific.
\\
\sphinxhline
\sphinxAtStartPar
☐
&
\sphinxAtStartPar
Abstract explains what, where, when, and why.
\\
\sphinxhline
\sphinxAtStartPar
☐
&
\sphinxAtStartPar
At least one Pelagos keyword is present (\sphinxcode{\sphinxupquote{pelagos sanctuary}} or
\sphinxcode{\sphinxupquote{pelagos agreement}}).
\\
\sphinxhline
\sphinxAtStartPar
☐
&
\sphinxAtStartPar
Category is set.
\\
\sphinxhline
\sphinxAtStartPar
☐
&
\sphinxAtStartPar
Spatial extent is defined and covers the correct area.
\\
\sphinxhline
\sphinxAtStartPar
☐
&
\sphinxAtStartPar
Date fields are filled in.
\\
\sphinxhline
\sphinxAtStartPar
☐
&
\sphinxAtStartPar
Licence is set.
\\
\sphinxhline
\sphinxAtStartPar
☐
&
\sphinxAtStartPar
Point of contact is identified.
\\
\sphinxhline
\sphinxAtStartPar
☐
&
\sphinxAtStartPar
Thumbnail is representative.
\\
\sphinxbottomrule
\end{tabular}
\sphinxtableafterendhook\par
\sphinxattableend\end{savenotes}

\sphinxstepscope


\chapter{Exploring Maps}
\label{\detokenize{04_exploring_maps:exploring-maps}}\label{\detokenize{04_exploring_maps:maps-en}}\label{\detokenize{04_exploring_maps::doc}}
\begin{sphinxcontents}
\sphinxstylecontentstitle{On this page}
\begin{itemize}
\item {} 
\sphinxAtStartPar
\phantomsection\label{\detokenize{04_exploring_maps:id1}}{\hyperref[\detokenize{04_exploring_maps:finding-a-map}]{\sphinxcrossref{Finding a Map}}}

\item {} 
\sphinxAtStartPar
\phantomsection\label{\detokenize{04_exploring_maps:id2}}{\hyperref[\detokenize{04_exploring_maps:the-map-viewer-interface}]{\sphinxcrossref{The Map Viewer Interface}}}

\item {} 
\sphinxAtStartPar
\phantomsection\label{\detokenize{04_exploring_maps:id3}}{\hyperref[\detokenize{04_exploring_maps:navigating-the-map}]{\sphinxcrossref{Navigating the Map}}}
\begin{itemize}
\item {} 
\sphinxAtStartPar
\phantomsection\label{\detokenize{04_exploring_maps:id4}}{\hyperref[\detokenize{04_exploring_maps:basic-navigation}]{\sphinxcrossref{Basic Navigation}}}

\item {} 
\sphinxAtStartPar
\phantomsection\label{\detokenize{04_exploring_maps:id5}}{\hyperref[\detokenize{04_exploring_maps:going-to-a-specific-location}]{\sphinxcrossref{Going to a Specific Location}}}

\end{itemize}

\item {} 
\sphinxAtStartPar
\phantomsection\label{\detokenize{04_exploring_maps:id6}}{\hyperref[\detokenize{04_exploring_maps:working-with-layers}]{\sphinxcrossref{Working with Layers}}}
\begin{itemize}
\item {} 
\sphinxAtStartPar
\phantomsection\label{\detokenize{04_exploring_maps:id7}}{\hyperref[\detokenize{04_exploring_maps:opening-the-layer-panel}]{\sphinxcrossref{Opening the Layer Panel}}}

\item {} 
\sphinxAtStartPar
\phantomsection\label{\detokenize{04_exploring_maps:id8}}{\hyperref[\detokenize{04_exploring_maps:toggling-layer-visibility}]{\sphinxcrossref{Toggling Layer Visibility}}}

\item {} 
\sphinxAtStartPar
\phantomsection\label{\detokenize{04_exploring_maps:id9}}{\hyperref[\detokenize{04_exploring_maps:changing-layer-order}]{\sphinxcrossref{Changing Layer Order}}}

\item {} 
\sphinxAtStartPar
\phantomsection\label{\detokenize{04_exploring_maps:id10}}{\hyperref[\detokenize{04_exploring_maps:adjusting-layer-opacity}]{\sphinxcrossref{Adjusting Layer Opacity}}}

\item {} 
\sphinxAtStartPar
\phantomsection\label{\detokenize{04_exploring_maps:id11}}{\hyperref[\detokenize{04_exploring_maps:viewing-the-legend}]{\sphinxcrossref{Viewing the Legend}}}

\end{itemize}

\item {} 
\sphinxAtStartPar
\phantomsection\label{\detokenize{04_exploring_maps:id12}}{\hyperref[\detokenize{04_exploring_maps:querying-features}]{\sphinxcrossref{Querying Features}}}
\begin{itemize}
\item {} 
\sphinxAtStartPar
\phantomsection\label{\detokenize{04_exploring_maps:id13}}{\hyperref[\detokenize{04_exploring_maps:clicking-on-a-feature}]{\sphinxcrossref{Clicking on a Feature}}}

\item {} 
\sphinxAtStartPar
\phantomsection\label{\detokenize{04_exploring_maps:id14}}{\hyperref[\detokenize{04_exploring_maps:identify-tool}]{\sphinxcrossref{Identify Tool}}}

\item {} 
\sphinxAtStartPar
\phantomsection\label{\detokenize{04_exploring_maps:id15}}{\hyperref[\detokenize{04_exploring_maps:selecting-and-filtering-features}]{\sphinxcrossref{Selecting and Filtering Features}}}

\end{itemize}

\item {} 
\sphinxAtStartPar
\phantomsection\label{\detokenize{04_exploring_maps:id16}}{\hyperref[\detokenize{04_exploring_maps:measurement-tools}]{\sphinxcrossref{Measurement Tools}}}
\begin{itemize}
\item {} 
\sphinxAtStartPar
\phantomsection\label{\detokenize{04_exploring_maps:id17}}{\hyperref[\detokenize{04_exploring_maps:distance-measurement}]{\sphinxcrossref{Distance Measurement}}}

\item {} 
\sphinxAtStartPar
\phantomsection\label{\detokenize{04_exploring_maps:id18}}{\hyperref[\detokenize{04_exploring_maps:area-measurement}]{\sphinxcrossref{Area Measurement}}}

\end{itemize}

\item {} 
\sphinxAtStartPar
\phantomsection\label{\detokenize{04_exploring_maps:id19}}{\hyperref[\detokenize{04_exploring_maps:printing-and-exporting-the-map-view}]{\sphinxcrossref{Printing and Exporting the Map View}}}
\begin{itemize}
\item {} 
\sphinxAtStartPar
\phantomsection\label{\detokenize{04_exploring_maps:id20}}{\hyperref[\detokenize{04_exploring_maps:print}]{\sphinxcrossref{Print}}}

\item {} 
\sphinxAtStartPar
\phantomsection\label{\detokenize{04_exploring_maps:id21}}{\hyperref[\detokenize{04_exploring_maps:share-permalink}]{\sphinxcrossref{Share / Permalink}}}

\item {} 
\sphinxAtStartPar
\phantomsection\label{\detokenize{04_exploring_maps:id22}}{\hyperref[\detokenize{04_exploring_maps:embed}]{\sphinxcrossref{Embed}}}

\end{itemize}

\item {} 
\sphinxAtStartPar
\phantomsection\label{\detokenize{04_exploring_maps:id23}}{\hyperref[\detokenize{04_exploring_maps:saving-changes-to-a-map}]{\sphinxcrossref{Saving Changes to a Map}}}

\item {} 
\sphinxAtStartPar
\phantomsection\label{\detokenize{04_exploring_maps:id24}}{\hyperref[\detokenize{04_exploring_maps:creating-a-new-map}]{\sphinxcrossref{Creating a New Map}}}

\end{itemize}
\end{sphinxcontents}

\sphinxAtStartPar
A \sphinxstylestrong{Map} in KMaP is an interactive, multi\sphinxhyphen{}layer web map built from one or more Datasets.
This section explains how to find, open, and navigate Maps, how to work with layers, and
how to extract information from the data displayed.


\section{Finding a Map}
\label{\detokenize{04_exploring_maps:finding-a-map}}
\sphinxAtStartPar
Maps are listed in the catalogue alongside all other resource types. To filter the
catalogue to show only Maps:
\begin{enumerate}
\sphinxsetlistlabels{\arabic}{enumi}{enumii}{}{.}%
\item {} 
\sphinxAtStartPar
Open the catalogue: click \sphinxstylestrong{Maps} in the top navigation bar.

\item {} 
\sphinxAtStartPar
In the filter panel on the left, under \sphinxstylestrong{Resource type}, select \sphinxstylestrong{Map}.

\item {} 
\sphinxAtStartPar
To find Pelagos\sphinxhyphen{}related maps, also apply the keyword filter:
in the \sphinxstylestrong{Keywords} field, type \sphinxcode{\sphinxupquote{pelagos sanctuary}} or \sphinxcode{\sphinxupquote{pelagos agreement}}.

\end{enumerate}

\sphinxAtStartPar
Alternatively, search for a map by title using the search bar.

\sphinxAtStartPar
Click on a Map card to open its \sphinxstylestrong{detail page}. From there, click \sphinxstylestrong{View map} to open
it in the interactive map viewer.


\section{The Map Viewer Interface}
\label{\detokenize{04_exploring_maps:the-map-viewer-interface}}
\sphinxAtStartPar
The map viewer (powered by \sphinxstylestrong{MapStore}) opens in a full\sphinxhyphen{}screen layout.
Its main interface elements are:


\begin{savenotes}\sphinxattablestart
\sphinxthistablewithglobalstyle
\centering
\begin{tabular}[t]{\X{30}{100}\X{70}{100}}
\sphinxtoprule
\sphinxstyletheadfamily 
\sphinxAtStartPar
Element
&\sphinxstyletheadfamily 
\sphinxAtStartPar
Description
\\
\sphinxmidrule
\sphinxtableatstartofbodyhook
\sphinxAtStartPar
\sphinxstylestrong{Main map canvas}
&
\sphinxAtStartPar
The central interactive map area.
\\
\sphinxhline
\sphinxAtStartPar
\sphinxstylestrong{Layer panel} (☰)
&
\sphinxAtStartPar
Toggles layer visibility, reorders layers, accesses layer settings.
\\
\sphinxhline
\sphinxAtStartPar
\sphinxstylestrong{Basemap switcher}
&
\sphinxAtStartPar
Changes the background map (OpenStreetMap, satellite imagery, …).
\\
\sphinxhline
\sphinxAtStartPar
\sphinxstylestrong{Zoom controls} (+ / −)
&
\sphinxAtStartPar
Zoom in and out. Also supports mouse wheel and pinch gesture.
\\
\sphinxhline
\sphinxAtStartPar
\sphinxstylestrong{Toolbar}
&
\sphinxAtStartPar
Row of tool icons at the top or side of the map canvas.
\\
\sphinxhline
\sphinxAtStartPar
\sphinxstylestrong{Scale bar}
&
\sphinxAtStartPar
Displays the current map scale.
\\
\sphinxhline
\sphinxAtStartPar
\sphinxstylestrong{Coordinates panel}
&
\sphinxAtStartPar
Shows the geographic coordinates of the mouse cursor position.
\\
\sphinxbottomrule
\end{tabular}
\sphinxtableafterendhook\par
\sphinxattableend\end{savenotes}


\section{Navigating the Map}
\label{\detokenize{04_exploring_maps:navigating-the-map}}

\subsection{Basic Navigation}
\label{\detokenize{04_exploring_maps:basic-navigation}}

\begin{savenotes}\sphinxattablestart
\sphinxthistablewithglobalstyle
\centering
\begin{tabular}[t]{\X{35}{100}\X{65}{100}}
\sphinxtoprule
\sphinxstyletheadfamily 
\sphinxAtStartPar
Action
&\sphinxstyletheadfamily 
\sphinxAtStartPar
How to perform it
\\
\sphinxmidrule
\sphinxtableatstartofbodyhook
\sphinxAtStartPar
\sphinxstylestrong{Pan (move the map)}
&
\sphinxAtStartPar
Click and drag on the map canvas.
\\
\sphinxhline
\sphinxAtStartPar
\sphinxstylestrong{Zoom in}
&
\sphinxAtStartPar
Scroll the mouse wheel forward, click \sphinxstylestrong{+}, or double\sphinxhyphen{}click on the map.
\\
\sphinxhline
\sphinxAtStartPar
\sphinxstylestrong{Zoom out}
&
\sphinxAtStartPar
Scroll the mouse wheel backward or click \sphinxstylestrong{−}.
\\
\sphinxhline
\sphinxAtStartPar
\sphinxstylestrong{Zoom to full extent}
&
\sphinxAtStartPar
Click the \sphinxstyleemphasis{Home} icon (🏠) in the toolbar to return to the default view.
\\
\sphinxhline
\sphinxAtStartPar
\sphinxstylestrong{Zoom to a layer}
&
\sphinxAtStartPar
In the Layer panel, right\sphinxhyphen{}click a layer and select \sphinxstyleemphasis{Zoom to layer extent}.
\\
\sphinxbottomrule
\end{tabular}
\sphinxtableafterendhook\par
\sphinxattableend\end{savenotes}


\subsection{Going to a Specific Location}
\label{\detokenize{04_exploring_maps:going-to-a-specific-location}}
\sphinxAtStartPar
Use the \sphinxstylestrong{Search / Geocoder} tool in the toolbar (🔍 icon) to search for a place name.
Type the name and press Enter — the map will zoom to that location.


\section{Working with Layers}
\label{\detokenize{04_exploring_maps:working-with-layers}}

\subsection{Opening the Layer Panel}
\label{\detokenize{04_exploring_maps:opening-the-layer-panel}}
\sphinxAtStartPar
Click the \sphinxstylestrong{☰} (hamburger) icon to open the layer panel. It lists all layers included
in the map, grouped into:
\begin{itemize}
\item {} 
\sphinxAtStartPar
\sphinxstylestrong{Overlays} – the thematic data layers (e.g. Pelagos Sanctuary boundary, cetacean
sighting points, shipping lanes).

\item {} 
\sphinxAtStartPar
\sphinxstylestrong{Background layers} – the basemap options (only one can be active at a time).

\end{itemize}


\subsection{Toggling Layer Visibility}
\label{\detokenize{04_exploring_maps:toggling-layer-visibility}}
\sphinxAtStartPar
Each layer in the panel has an \sphinxstylestrong{eye icon} (👁). Click it to show or hide that layer.
Hidden layers remain in the map composition but are not displayed on the canvas.


\subsection{Changing Layer Order}
\label{\detokenize{04_exploring_maps:changing-layer-order}}
\sphinxAtStartPar
Layers are drawn in the order they appear in the panel (top layers are drawn on top).
To reorder:
\begin{enumerate}
\sphinxsetlistlabels{\arabic}{enumi}{enumii}{}{.}%
\item {} 
\sphinxAtStartPar
Click and hold a layer name in the panel.

\item {} 
\sphinxAtStartPar
Drag it up or down to the desired position.

\item {} 
\sphinxAtStartPar
Release to drop.

\end{enumerate}


\subsection{Adjusting Layer Opacity}
\label{\detokenize{04_exploring_maps:adjusting-layer-opacity}}\begin{enumerate}
\sphinxsetlistlabels{\arabic}{enumi}{enumii}{}{.}%
\item {} 
\sphinxAtStartPar
In the Layer panel, click the \sphinxstylestrong{⋮} (options) icon next to a layer.

\item {} 
\sphinxAtStartPar
Select \sphinxstylestrong{Settings} (or the slider icon).

\item {} 
\sphinxAtStartPar
Use the \sphinxstylestrong{Opacity} slider to set transparency between 0 \% (invisible) and 100 \% (solid).

\end{enumerate}


\subsection{Viewing the Legend}
\label{\detokenize{04_exploring_maps:viewing-the-legend}}\begin{enumerate}
\sphinxsetlistlabels{\arabic}{enumi}{enumii}{}{.}%
\item {} 
\sphinxAtStartPar
Click the \sphinxstylestrong{⋮} (options) icon next to a layer.

\item {} 
\sphinxAtStartPar
Select \sphinxstylestrong{Legend} to expand the symbol key for that layer.

\end{enumerate}

\sphinxAtStartPar
Alternatively, some map configurations display a persistent legend panel on the side
of the canvas.


\section{Querying Features}
\label{\detokenize{04_exploring_maps:querying-features}}

\subsection{Clicking on a Feature}
\label{\detokenize{04_exploring_maps:clicking-on-a-feature}}
\sphinxAtStartPar
Click on any visible feature (point, line, or polygon) on the map canvas to open an
\sphinxstylestrong{attribute pop\sphinxhyphen{}up}. The pop\sphinxhyphen{}up displays the attribute values associated with that
feature — for example, the name of a marine protected area, a species name, or a
monitoring value.

\sphinxAtStartPar
If multiple layers overlap at the clicked location, the pop\sphinxhyphen{}up may show results from
all overlapping layers. Use the arrow buttons in the pop\sphinxhyphen{}up to navigate between them.


\subsection{Identify Tool}
\label{\detokenize{04_exploring_maps:identify-tool}}
\sphinxAtStartPar
The \sphinxstylestrong{Identify} tool (ℹ️ icon in the toolbar) works similarly to clicking directly
on the map. It queries all visible layers at the clicked point and returns their
attributes in a panel.


\subsection{Selecting and Filtering Features}
\label{\detokenize{04_exploring_maps:selecting-and-filtering-features}}
\sphinxAtStartPar
For vector layers, MapStore offers a \sphinxstylestrong{Filter} tool:
\begin{enumerate}
\sphinxsetlistlabels{\arabic}{enumi}{enumii}{}{.}%
\item {} 
\sphinxAtStartPar
In the Layer panel, click \sphinxstylestrong{⋮} next to the desired layer.

\item {} 
\sphinxAtStartPar
Select \sphinxstylestrong{Filter}.

\item {} 
\sphinxAtStartPar
Define a filter condition using the attribute fields
(e.g. show only records where \sphinxstyleemphasis{country = “France”}).

\item {} 
\sphinxAtStartPar
Click \sphinxstylestrong{Apply}. Only features matching the filter are displayed.

\end{enumerate}

\begin{sphinxadmonition}{tip}{Tip:}
\sphinxAtStartPar
Use the filter tool to focus on a specific subset of Pelagos data — for example,
showing only cetacean sightings from a specific year, or only routes that pass
through the sanctuary.
\end{sphinxadmonition}


\section{Measurement Tools}
\label{\detokenize{04_exploring_maps:measurement-tools}}

\subsection{Distance Measurement}
\label{\detokenize{04_exploring_maps:distance-measurement}}\begin{enumerate}
\sphinxsetlistlabels{\arabic}{enumi}{enumii}{}{.}%
\item {} 
\sphinxAtStartPar
Click the \sphinxstylestrong{Measure} tool (📐 icon) in the toolbar.

\item {} 
\sphinxAtStartPar
Select \sphinxstylestrong{Distance}.

\item {} 
\sphinxAtStartPar
Click points on the map to define a line. Double\sphinxhyphen{}click to end.

\item {} 
\sphinxAtStartPar
The total distance is displayed in the unit of your choice (km, miles, nm).

\end{enumerate}


\subsection{Area Measurement}
\label{\detokenize{04_exploring_maps:area-measurement}}\begin{enumerate}
\sphinxsetlistlabels{\arabic}{enumi}{enumii}{}{.}%
\item {} 
\sphinxAtStartPar
Click the \sphinxstylestrong{Measure} tool.

\item {} 
\sphinxAtStartPar
Select \sphinxstylestrong{Area}.

\item {} 
\sphinxAtStartPar
Click to define the vertices of a polygon. Double\sphinxhyphen{}click to close.

\item {} 
\sphinxAtStartPar
The calculated area is displayed.

\end{enumerate}

\begin{sphinxadmonition}{note}{Note:}
\sphinxAtStartPar
Measurements are performed on the map projection and may have small geometric
inaccuracies for large areas. For precise measurements, use a desktop GIS application.
\end{sphinxadmonition}


\section{Printing and Exporting the Map View}
\label{\detokenize{04_exploring_maps:printing-and-exporting-the-map-view}}

\subsection{Print}
\label{\detokenize{04_exploring_maps:print}}\begin{enumerate}
\sphinxsetlistlabels{\arabic}{enumi}{enumii}{}{.}%
\item {} 
\sphinxAtStartPar
Click the \sphinxstylestrong{Print} tool (🖨 icon) in the toolbar.

\item {} 
\sphinxAtStartPar
A print dialog appears. Configure:
\begin{itemize}
\item {} 
\sphinxAtStartPar
\sphinxstylestrong{Title} – a title to display on the printed map.

\item {} 
\sphinxAtStartPar
\sphinxstylestrong{Paper size} – A4, A3, Letter, etc.

\item {} 
\sphinxAtStartPar
\sphinxstylestrong{Resolution} – 72, 150, or 300 dpi.

\item {} 
\sphinxAtStartPar
\sphinxstylestrong{Layout} – portrait or landscape.

\end{itemize}

\item {} 
\sphinxAtStartPar
Click \sphinxstylestrong{Print} to generate a PDF or image.

\end{enumerate}


\subsection{Share / Permalink}
\label{\detokenize{04_exploring_maps:share-permalink}}
\sphinxAtStartPar
To share the current map view (including zoom level, active layers, and centre point):
\begin{enumerate}
\sphinxsetlistlabels{\arabic}{enumi}{enumii}{}{.}%
\item {} 
\sphinxAtStartPar
Click the \sphinxstylestrong{Share} button (link icon) in the toolbar.

\item {} 
\sphinxAtStartPar
A permalink URL is generated. Copy and share this URL — recipients will open the map
in the same state.

\end{enumerate}


\subsection{Embed}
\label{\detokenize{04_exploring_maps:embed}}\begin{enumerate}
\sphinxsetlistlabels{\arabic}{enumi}{enumii}{}{.}%
\item {} 
\sphinxAtStartPar
Click the \sphinxstylestrong{Share} button.

\item {} 
\sphinxAtStartPar
Switch to the \sphinxstylestrong{Embed} tab.

\item {} 
\sphinxAtStartPar
Copy the HTML \sphinxcode{\sphinxupquote{<iframe>}} code and paste it into your website or report.

\end{enumerate}


\section{Saving Changes to a Map}
\label{\detokenize{04_exploring_maps:saving-changes-to-a-map}}
\sphinxAtStartPar
If you have \sphinxstylestrong{edit permission} on a Map, you can save changes to layers, styles, and
the initial view:
\begin{enumerate}
\sphinxsetlistlabels{\arabic}{enumi}{enumii}{}{.}%
\item {} 
\sphinxAtStartPar
Make your changes in the map viewer (reorder layers, change styles, set a new
initial extent, etc.).

\item {} 
\sphinxAtStartPar
Click the \sphinxstylestrong{Save} button (💾 icon) in the toolbar.

\item {} 
\sphinxAtStartPar
Choose \sphinxstylestrong{Save} (overwrite the existing map) or \sphinxstylestrong{Save as} (create a new map).

\end{enumerate}

\begin{sphinxadmonition}{warning}{Warning:}
\sphinxAtStartPar
Saving overwrites the map for all users who have access to it.
If you want to experiment without affecting others, use \sphinxstylestrong{Save as} to create a
personal copy first.
\end{sphinxadmonition}


\section{Creating a New Map}
\label{\detokenize{04_exploring_maps:creating-a-new-map}}
\sphinxAtStartPar
Members of the Pelagos Agreement group with appropriate permissions can create new Maps:
\begin{enumerate}
\sphinxsetlistlabels{\arabic}{enumi}{enumii}{}{.}%
\item {} 
\sphinxAtStartPar
Go to \sphinxstylestrong{Maps} in the navigation bar.

\item {} 
\sphinxAtStartPar
Click \sphinxstylestrong{+ Create map} (or the equivalent button at the top of the catalogue).

\item {} 
\sphinxAtStartPar
The map viewer opens with an empty canvas.

\item {} 
\sphinxAtStartPar
Use the \sphinxstylestrong{Add layer} button (+ icon in the Layer panel) to search for and add
Datasets from the KMaP catalogue.

\item {} 
\sphinxAtStartPar
Style each layer as needed.

\item {} 
\sphinxAtStartPar
Click \sphinxstylestrong{Save} and provide a title, abstract, and keywords
(remember to include \sphinxcode{\sphinxupquote{pelagos sanctuary}} or \sphinxcode{\sphinxupquote{pelagos agreement}}).

\end{enumerate}



\renewcommand{\indexname}{Index}
\printindex
\end{document}